\documentclass[12pt]{article}
\usepackage[margin=1in]{geometry}
\usepackage{amsmath,amssymb,amsthm,booktabs,longtable,array,enumitem,float}
\usepackage{newtxtext,newtxmath}
\usepackage[hidelinks]{hyperref}
% p5_numbers.tex: generated by p5_observables.py from the P5 registry rows.
% Prints with \reg{<key>}.  Do not edit by hand.
\makeatletter
\providecommand{\reg}[1]{\ifcsname reg@#1\endcsname\csname reg@#1\endcsname\else\textbf{??\detokenize{#1}}\fi}
\expandafter\def\csname reg@u_p5_B\endcsname{1.718282}
\expandafter\def\csname reg@lambda_p5_B\endcsname{0.336122}
\expandafter\def\csname reg@beta_p5_B\endcsname{0.294655}
\expandafter\def\csname reg@b_p5_B\endcsname{0.195615}
\expandafter\def\csname reg@M_p5_B\endcsname{0.731059}
\expandafter\def\csname reg@m_p5_B\endcsname{0.268941}
\expandafter\def\csname reg@d_slot_B\endcsname{0.363676}
\expandafter\def\csname reg@vol_none_B\endcsname{0.391230}
\expandafter\def\csname reg@vol_buy_B\endcsname{0.559291}
\expandafter\def\csname reg@vol_full_B\endcsname{0.727352}
\expandafter\def\csname reg@theta_B\endcsname{0.462117}
\expandafter\def\csname reg@theta_pct_B\endcsname{46.2\%}
\expandafter\def\csname reg@margin_obs_beta_B\endcsname{0.038678}
\expandafter\def\csname reg@margin_obs_u_hi_B\endcsname{0.004617}
\expandafter\def\csname reg@margin_obs_buy_recruit_B\endcsname{0.026059}
\expandafter\def\csname reg@margin_obs_ceiling_B\endcsname{0.000334}
\expandafter\def\csname reg@rC_check_B\endcsname{3.592659}
\expandafter\def\csname reg@margin_obs_r2_above_rC_B\endcsname{0.007341}
\expandafter\def\csname reg@margin_obs_full_orders_r1_B\endcsname{0.024824}
\expandafter\def\csname reg@margin_obs_full_orders_r2_B\endcsname{0.043127}
\expandafter\def\csname reg@margin_obs_weak_r0_B\endcsname{0.003333}
\expandafter\def\csname reg@margin_obs_neutral_r1_B\endcsname{1.708333}
\expandafter\def\csname reg@rF_check_B\endcsname{1.467127}
\expandafter\def\csname reg@rS_p5_B\endcsname{2.124148}
\expandafter\def\csname reg@rN_check_B\endcsname{1.747878}
\expandafter\def\csname reg@margin_obs_rF_B\endcsname{0.082873}
\expandafter\def\csname reg@margin_obs_rN_B\endcsname{0.097878}
\expandafter\def\csname reg@margin_obs_lowsell_r165_B\endcsname{0.003996}
\expandafter\def\csname reg@margin_obs_buy_r155_B\endcsname{0.006243}
\expandafter\def\csname reg@margin_obs_buy_recruit_r165_B\endcsname{0.713202}
\expandafter\def\csname reg@theta_lo\endcsname{0.410000}
\expandafter\def\csname reg@theta_hi\endcsname{0.462784}
\expandafter\def\csname reg@own_p5_r1_B\endcsname{0.324401}
\expandafter\def\csname reg@entry_floor_obs\endcsname{0.25}
\expandafter\def\csname reg@entry_plateau_B\endcsname{0.522757}
\expandafter\def\csname reg@entry_jump_B\endcsname{0.272757}
\expandafter\def\csname reg@entry_r2_obs_B\endcsname{0.25}
\expandafter\def\csname reg@ba_class_r1_B\endcsname{0.077020}
\expandafter\def\csname reg@ba_entry_r1_B\endcsname{0.709279}
\expandafter\def\csname reg@ba_r1_B\endcsname{0.786299}
\expandafter\def\csname reg@ba_share_r1_B\endcsname{90.2\%}
\expandafter\def\csname reg@ba_ratio_r1_B\endcsname{10.2}
\expandafter\def\csname reg@impact_sell_r1_B\endcsname{0.038510}
\expandafter\def\csname reg@impact_buy_r1_B\endcsname{0.747789}
\expandafter\def\csname reg@impact_asym_r1_B\endcsname{0.709279}
\expandafter\def\csname reg@ba_r2_B\endcsname{0.108469}
\expandafter\def\csname reg@impact_sell_r2_B\endcsname{0.054235}
\expandafter\def\csname reg@pbar_r1_B\endcsname{0.872531}
\expandafter\def\csname reg@p_neutral_r1_B\endcsname{0.614583}
\expandafter\def\csname reg@pdrop_neutral_r1_B\endcsname{-0.257948}
\expandafter\def\csname reg@pdrop_neutral_r2_B\endcsname{0.00}
\expandafter\def\csname reg@ba_rC_left_B\endcsname{0.969913}
\expandafter\def\csname reg@ba_rC_right_B\endcsname{0.108078}
\expandafter\def\csname reg@ba_jump_rC_B\endcsname{0.861835}
\expandafter\def\csname reg@pbar_jump_rC_B\endcsname{-0.313429}
\expandafter\def\csname reg@alt_lambda_B\endcsname{0.323267}
\expandafter\def\csname reg@alt_beta_B\endcsname{0.298555}
\expandafter\def\csname reg@alt_theta_B\endcsname{0.444444}
\expandafter\def\csname reg@alt_ba_r1_B\endcsname{0.778935}
\expandafter\def\csname reg@alt_entry_r1_B\endcsname{0.522757}
\expandafter\def\csname reg@theta_min_counts_B\endcsname{0.250299}
\expandafter\def\csname reg@imb_buy_B\endcsname{0.168061}
\expandafter\def\csname reg@theta_buy_B\endcsname{0.300489}
\expandafter\def\csname reg@theta_buy_pct_B\endcsname{30.0\%}
\expandafter\def\csname reg@U_entry_B\endcsname{0.522757}
\expandafter\def\csname reg@ba_U_r160_B\endcsname{0.384588}
\expandafter\def\csname reg@U_gap_r160_B\endcsname{0.005363}
\expandafter\def\csname reg@U_mu0_B\endcsname{0.309329}
\expandafter\def\csname reg@U_sell_mu_B\endcsname{0.50}
\expandafter\def\csname reg@ba_U_r155_B\endcsname{0.373656}
\expandafter\def\csname reg@ba_U_r165_B\endcsname{0.395776}
\expandafter\def\csname reg@wgap_r1_B\endcsname{0.308078}
\expandafter\def\csname reg@wgap_r2_B\endcsname{0.333751}
\expandafter\def\csname reg@wprep_buy_r1_B\endcsname{0.820706}
\expandafter\def\csname reg@wprep_sell_r1_B\endcsname{0.512628}
\expandafter\def\csname reg@wprep_neutral_r1_B\endcsname{0.666667}
\expandafter\def\csname reg@wprep_buy_r2_B\endcsname{0.805765}
\expandafter\def\csname reg@wprep_sell_r2_B\endcsname{0.472013}
\expandafter\def\csname reg@wprep_neutral_r2_B\endcsname{0.638889}
\expandafter\def\csname reg@cov_entry_r1_B\endcsname{0.272757}
\expandafter\def\csname reg@cov_entry_r2_B\endcsname{0.00}
\expandafter\def\csname reg@cov_win_r1_B\endcsname{0.251863}
\expandafter\def\csname reg@cov_win_r2_B\endcsname{0.030344}
\expandafter\def\csname reg@win_prob_r1_B\endcsname{0.390520}
\expandafter\def\csname reg@win_prob_r2_B\endcsname{0.159722}
\expandafter\def\csname reg@evA_r1_sale_B\endcsname{0.478221}
\expandafter\def\csname reg@evAr_r1_sale_B\endcsname{0.830138}
\expandafter\def\csname reg@evA_r1_neutral_B\endcsname{0.593750}
\expandafter\def\csname reg@evAr_r1_neutral_B\endcsname{0.966102}
\expandafter\def\csname reg@evA_r1_buy_B\endcsname{0.000000}
\expandafter\def\csname reg@evN_r1_sale_B\endcsname{-0.159407}
\expandafter\def\csname reg@evNr_r1_sale_B\endcsname{-0.276713}
\expandafter\def\csname reg@evN_r1_neutral_B\endcsname{-0.197917}
\expandafter\def\csname reg@evNr_r1_neutral_B\endcsname{-0.322034}
\expandafter\def\csname reg@evA_r2_sale_B\endcsname{0.538338}
\expandafter\def\csname reg@evAr_r2_sale_B\endcsname{0.882519}
\expandafter\def\csname reg@evA_r2_neutral_B\endcsname{0.701042}
\expandafter\def\csname reg@evAr_r2_neutral_B\endcsname{1.055410}
\expandafter\def\csname reg@evA_r2_buy_B\endcsname{0.863745}
\expandafter\def\csname reg@evAr_r2_buy_B\endcsname{1.202200}
\expandafter\def\csname reg@evN_r2_sale_B\endcsname{-0.179446}
\expandafter\def\csname reg@evNr_r2_sale_B\endcsname{-0.294173}
\expandafter\def\csname reg@evN_r2_neutral_B\endcsname{-0.233681}
\expandafter\def\csname reg@evNr_r2_neutral_B\endcsname{-0.351803}
\expandafter\def\csname reg@evN_r2_buy_B\endcsname{-0.287915}
\expandafter\def\csname reg@evNr_r2_buy_B\endcsname{-0.400733}
\expandafter\def\csname reg@price_r1_sale_B\endcsname{0.576074}
\expandafter\def\csname reg@price_r1_neutral_B\endcsname{0.614583}
\expandafter\def\csname reg@price_r1_buy_B\endcsname{1.362372}
\expandafter\def\csname reg@price_r2_sale_B\endcsname{0.610002}
\expandafter\def\csname reg@price_r2_neutral_B\endcsname{0.664236}
\expandafter\def\csname reg@price_r2_buy_B\endcsname{0.718471}
\expandafter\def\csname reg@delta_rS_B\endcsname{0.297463}
\expandafter\def\csname reg@rS_cert_hi_B\endcsname{2.124149}
\expandafter\def\csname reg@rS_counter_r_B\endcsname{2.1241481}
\expandafter\def\csname reg@rS_counter_margin_B\endcsname{\ensuremath{-8.59\times10^{-10}}}
\expandafter\def\csname reg@theta_max_full_eq_r1_B\endcsname{0.760000}
\expandafter\def\csname reg@theta_max_full_eq_r2_B\endcsname{0.829586}
\expandafter\def\csname reg@margin_obs_band_incentive_B\endcsname{0.024768}
\expandafter\def\csname reg@theta09_beta_B\endcsname{0.105296}
\expandafter\def\csname reg@theta09_lowsale_r1_B\endcsname{-0.011667}
\expandafter\def\csname reg@margin_obs_u_lo_B\endcsname{0.328451}
\expandafter\def\csname reg@margin_obs_post_over_58_B\endcsname{0.106059}
\expandafter\def\csname reg@margin_obs_rS_rC_B\endcsname{1.468510}
\expandafter\def\csname reg@margin_obs_rN_rC_B\endcsname{1.844781}
\expandafter\def\csname reg@margin_obs_rS_h_B\endcsname{7.875852}
\expandafter\def\csname reg@mix_sigma_r160_B\endcsname{0.492845}
\expandafter\def\csname reg@mix_e1_r160_B\endcsname{0.506106}
\expandafter\def\csname reg@mix_lowsale_r160_B\endcsname{-0.005938}
\expandafter\def\csname reg@mix_mean_order_r160_B\endcsname{0.082828}
\expandafter\def\csname reg@mix_entry_r160_B\endcsname{0.321311}
\expandafter\def\csname reg@evprep_mean_r1_B\endcsname{0.160592}
\expandafter\def\csname reg@evprep_meanr_r1_B\endcsname{0.270349}
\expandafter\def\csname reg@evprep_mean_r2_B\endcsname{0.701042}
\expandafter\def\csname reg@evprep_meanr_r2_B\endcsname{1.045918}
\makeatother


\newtheorem{lemma}{Lemma}
\newtheorem{proposition}{Proposition}
\theoremstyle{definition}
\newtheorem{definition}{Definition}
\newtheorem{remark}{Remark}

\newcommand{\E}{\mathbb E}
\newcommand{\Ent}{\mathsf E}
\newcommand{\fr}{\mathfrak r}
\newcommand{\one}{\mathbf 1}
\newcommand{\cH}{c_H}
\newcommand{\status}[1]{\textsc{#1}}
\DeclareMathOperator{\Cov}{Cov}
\setlist{nosep}

\title{Empirical implications of the order slot\\
\large Proof package P5 (result R11)}
\author{}
\date{23 October 2026 (round 2)}

\begin{document}
\maketitle

\section*{Summary}

This package derives the empirical implications of the order slot.
Each statement has a proof and one status label: analytical, computer-assisted, or input.
Values at calibration B come from interval arithmetic.
The script is \texttt{scripts/P5\_observables/p5\_observables.py}.
Its run log is \texttt{scripts/P5\_observables/run\_log.txt}.
The registry rows are in \texttt{registry/P5\_observables.csv}.

Table~\ref{tab:summary} lists the results.
Section~\ref{sec:scope} states the scope and lists what stays open.
Section~\ref{sec:old} compares each result with the earlier study.
Section~\ref{sec:margins} lists the margins, the tie rules, and the parameter set.

\begin{table}[htbp]
\centering\small
\caption{Results of the package.}\label{tab:summary}
\begin{tabular}{@{}p{0.9cm}p{6.6cm}p{2.3cm}p{4.6cm}@{}}
\toprule
No. & Statement & Status & Value at B\\
\midrule
Prop.~\ref{p:flow} & Volume, mean order, and informed share identify the outcome: no informed trade, buys only, or buys and sales. & analytical; values computer-assisted & volume \reg{vol_none_B}, \reg{vol_buy_B}, \reg{vol_full_B}; informed share \reg{theta_buy_B}, \reg{theta_B}\\
Prop.~\ref{p:spread} & The conditional execution spread is a class component plus an entry component. Buy impact exceeds sell impact exactly by the entry component. The spread falls at the ceiling $r_C$. Above $r_C$ the entry component is zero. & analytical; values computer-assisted & $S(r_1)=\reg{ba_r1_B}$; entry share \reg{ba_share_r1_B}; fall at $r_C=\reg{rC_check_B}$: \reg{ba_jump_rC_B}\\
Prop.~\ref{p:ident} & Order counts leave the informed share open on a range that the investor incentives bound. Sell impact and volume identify it on the two-sided branch. Counts identify it on the buy-only branch. & analytical; values computer-assisted & counts and incentives at $r_1$: $[\reg{theta_min_counts_B},\reg{theta_max_full_eq_r1_B})$; band for $\vartheta$: (\reg{theta_lo}, \reg{theta_hi})\\
Prop.~\ref{p:entry} & Entry is the floor plus $(1-\rho)$ times the buy frequency, below $r_C$. Entry is flat at the floor above $r_C$. A hidden price gives a flat entry. & analytical; values computer-assisted & plateau \reg{entry_plateau_B}; floor \reg{entry_floor_obs}\\
Prop.~\ref{p:sec7} & The Section 7 predictions of the Laplace manuscript hold in the slot. Strength raises entry when the price is visible. Strong competition stops the entry response and keeps the information. & analytical; values computer-assisted & entry rise \reg{entry_jump_B}; entry \reg{entry_r2_obs_B} at $r_2$\\
Prop.~\ref{p:win} & Orders predict the winner when they no longer predict entry. The winner gap is $\vartheta(1-\ell/r)$. & analytical; values computer-assisted & gap \reg{wgap_r1_B} at $r_1$, \reg{wgap_r2_B} at $r_2$\\
Prop.~\ref{p:event} & Announcement returns are exact. A neutral order lowers the price below $r_C$. The mean return over both announcements is zero. The mean over preparation announcements alone is positive. & analytical; values computer-assisted & neutral price fall \reg{pdrop_neutral_r1_B}; preparation sample \reg{evprep_mean_r1_B} at $r_1$\\
Prop.~\ref{p:coex} & On the buy-only branch the net order flow is positive and a sale executes above a no-trade price. Under the challenger rule, $\mathcal N$ and $\mathcal U$ are the only outcomes. & analytical; values computer-assisted & mean order \reg{imb_buy_B}; price gap \reg{U_gap_r160_B}\\
\bottomrule
\end{tabular}
\end{table}

\section{Set-up}\label{sec:setup}

\subsection{Primitives and notation}

The trading stage is the order slot of the brief.
The informed investor fills the slot with probability $\lambda\in(0,1)$ and chooses $q\in\{-1,0,1\}$.
A nonzero order costs $k>0$.
A liquidity trader fills the slot with probability $1-\lambda$.
It sells, abstains, or buys one unit with probabilities $\beta$, $1-2\beta$, $\beta$, where $0<\beta\le 1/3$.
Dealers see the order $X$ and set $P=\E[V_T\mid X]$.
The challenger sees the price and then decides whether to prepare.
This package fixes the prior $\pi=1/2$.

Write
\[
b=(1-\lambda)\beta,\qquad u=\frac{\lambda}{b},\qquad d=b+\frac{\lambda}{2},\qquad
M=\frac{1+u}{2+u},\qquad m=1-M,\qquad \vartheta=\frac{u}{2+u}=2M-1 .
\]
Here $M=M(1/2)$ and $m=m(1/2)$ are the posterior bounds of the brief.
The number $\vartheta=\lambda/(\lambda+2b)$ is the informed share of nonzero orders when both classes trade fully.

The auction primitives are $h=10$, $\ell=1$, $p=1/2$, $\rho=1/4$, $c_L=0$, $c_H=6$, $k=1/50$.
The incumbent value is uniform on $[0,r]$ with $\ell<r<h$, so $1<r<10$ at B.
The strengths are $r_0=6/5$, $r_1=3$, $r_2=18/5$.
For $p<\ell<r<h$ the shared objects are
\[
\begin{aligned}
t_0&=p\Bigl(1-\frac pr\Bigr), & t_H&=\frac r2+\frac{p^2}{2r}, & t_L&=\ell-\frac{\ell^2-p^2}{2r},\\
g_H&=h-t_H, & g_L&=\frac{\ell^2-p^2}{2r}, & \Delta_T&=t_H-t_L=\frac{(r-\ell)^2}{2r}.
\end{aligned}
\]
These forms equal equation (4) of the pinned manuscript (commit \texttt{01bd5d6}).
Section~0 of the run log integrates over $R$ at $r=3$ and $r=18/5$.
It reproduces $t_0,t_H,t_L,g_H,g_L,\Delta_T$ in exact rational arithmetic.

Define, for a belief $\mu\in[0,1]$,
\[
B_r(\mu)=g_L+\mu(g_H-g_L),\qquad
A_r(\mu)=t_L-t_0+\Delta_T(r)\,\mu,\qquad
\tau(r)=\frac{\cH-g_L}{g_H-g_L}.
\]
The challenger with cost $c_H$ prepares if and only if $B_r(\mu)\ge\cH$, that is $\mu\ge\tau(r)$.
The entry rule is $e_r(\mu)=\rho+(1-\rho)\one\{B_r(\mu)\ge\cH\}$.
Let $J_r=\one\{B_r(M)\ge\cH\}$.
Let $\fr(D)=1+D+\sqrt{D^2+2D}$.
For $\ell=1$ and $r>1$, the equation $\Delta_T(r)=D$ has the unique solution $r=\fr(D)$.
Define three strengths:
\[
r_N=\fr\Bigl(\frac{2k}{\rho}\Bigr),\qquad
r_F=\fr\Bigl(\frac{k}{m}\Bigr),\qquad
r_S=\fr\Bigl(\frac{k}{\rho\, m}\Bigr).
\]
The strength $r_S$ is the unique root $r>1$ of $\rho\,\Delta_T(r)\,m=k$.
At B, $m=1/(1+e)$, so $k/(\rho m)=2(1+e)/25=\reg{delta_rS_B}$ and $r_S=\reg{rS_p5_B}$ (rounded).
The exact $r_S$ is the endpoint in every statement below.
The rounded value $\reg{rS_p5_B}$ lies below $r_S$ and is not an endpoint.
The rounded value $\reg{rS_cert_hi_B}$ lies above $r_S$ and is a sufficient endpoint.

\begin{lemma}[Auction objects]\label{l:auction}
Let $p=1/2$, $\ell=1$, $h=10$, $c_H=6$ and $1<r<h$.
\begin{enumerate}[label=(\alph*)]
\item $t_L-t_0=\tfrac12-\tfrac1{8r}>0$. The function $A_r(\mu)$ is positive for $\mu\ge 0$, strictly increasing in $\mu$, and strictly increasing in $r$. Also $\Delta_T$ is strictly increasing in $r$.
\item $B_r(\tfrac12)=5-\tfrac r4+\tfrac1{8r}<\cH$.
\item $B_r(\mu)$ is strictly increasing in $\mu$ and strictly decreasing in $r$.
\item If $M>5/8$, the equation $B_r(M)=\cH$ has exactly one root $r_C(M)\in(1,h)$.
It is
\[
r_C(M)=\frac{(10M-6)+\sqrt{(10M-6)^2+M(\tfrac34-M)}}{M}.
\]
Then $J_r=1$ if and only if $r\le r_C(M)$.
\end{enumerate}
\end{lemma}

\begin{proof}
(a) Direct substitution gives $t_L-t_0=(\ell-p)-(\ell^2-3p^2)/(2r)=\tfrac12-\tfrac1{8r}$.
This is positive for $r>1/4$ and increases in $r$ with slope $1/(8r^2)$.
Also $\Delta_T'(r)=\tfrac12-\tfrac1{2r^2}>0$ for $r>1$.
Since $A_r(\mu)=(t_L-t_0)+\Delta_T\mu$, both terms are nonnegative and increase in $r$.
The slope in $\mu$ is $\Delta_T>0$.

(b) $B_r(\tfrac12)=\tfrac12(g_H+g_L)=\tfrac12\bigl[h-\tfrac r2+\tfrac{\ell^2-2p^2}{2r}\bigr]=5-\tfrac r4+\tfrac1{8r}$.
For $r>1$ this is below $5-\tfrac14+\tfrac18<6$.

(c) The slope in $\mu$ is $g_H-g_L=h-\tfrac r2-\tfrac{\ell^2}{2r}>0$ for $r<10$.
Also $g_H'(r)=-\tfrac12+\tfrac{p^2}{2r^2}<0$ and $g_L'(r)=-\tfrac{\ell^2-p^2}{2r^2}<0$.
Hence $B_r(\mu)$ falls in $r$ for each $\mu\ge 0$.

(d) One computes $B_r(M)=Mh-\tfrac{Mr}2+\tfrac{(1-M)\ell^2-p^2}{2r}$.
Its slope in $r$ is $-\tfrac{Mr^2+(1-M)\ell^2-p^2}{2r^2}$.
For $r>\ell$ the numerator exceeds $\ell^2-p^2>0$, so $B_r(M)$ falls strictly.
At $r=1$ it equals $9M+\tfrac38$, which exceeds $6$ if and only if $M>5/8$.
At $r=10$ it equals $4.95M+0.0375<6$.
So the root exists and is unique.
Multiply $B_r(M)=\cH$ by $2r$ to get $Mr^2-2(Mh-\cH)r-[(1-M)\ell^2-p^2]=0$.
The positive root is the stated formula.
\end{proof}

\subsection{Beliefs, prices and gains}

An informed profile is a pair $(\sigma_H,\sigma_L)$ of distributions on $\{-1,0,1\}$.
Write $a_j(x)=(1-\lambda)\eta(x)+\lambda\sigma_j(x)$, where $\eta=(\beta,1-2\beta,\beta)$ on $(-1,0,1)$.
Every order has positive liquidity mass, so Bayes' rule defines $\mu_x=\pi a_H(x)/[\pi a_H(x)+(1-\pi)a_L(x)]$ at every order.
No off-path belief is needed.

\begin{lemma}[Posterior bounds]\label{l:bounds}
For every informed profile and every $x\in\{-1,0,1\}$, $m\le\mu_x\le M$.
\end{lemma}

\begin{proof}
Put $c_x=(1-\lambda)\eta(x)$.
Then $c_{\pm1}=b$ and $c_0=(1-\lambda)(1-2\beta)\ge b$, because $\beta\le 1/3$.
Since $\sigma_j(x)\in[0,1]$, we have $c_x\le a_j(x)\le c_x+\lambda$.
Therefore
\[
\frac{a_H(x)}{a_L(x)}\le\frac{c_x+\lambda}{c_x}=1+\frac{\lambda}{c_x}\le 1+\frac\lambda b=1+u .
\]
The same argument gives $a_H/a_L\ge 1/(1+u)$.
The posterior at $\pi=\tfrac12$ equals $a_H/(a_H+a_L)$, and it increases in $a_H/a_L$.
The two bounds give $M=(1+u)/(2+u)$ and $m=1/(2+u)=1-M$.
\end{proof}

\begin{lemma}[Prices and gains]\label{l:prices}
Fix a candidate equilibrium.
Let $e_x\in[\rho,1]$ be the probability that the challenger prepares after the price $P(x)$ of the order $x$.
\begin{enumerate}[label=(\alph*)]
\item $\E[V_T\mid j,x]=t_0+e_x(t_j-t_0)$ for $j\in\{H,L\}$.
\item $P(x)=t_0+e_xA_r(\mu_x)$.
\item $\E[V_T\mid H,x]-P(x)=e_x\Delta_T(1-\mu_x)$ and $P(x)-\E[V_T\mid L,x]=e_x\Delta_T\mu_x$.
\item An informed investor of class $j$ who places $q$ at the fixed price and entry schedule earns
$U_H(q)=q\,e_q\Delta_T(1-\mu_q)-k|q|$ and $U_L(q)=-q\,e_q\Delta_T\mu_q-k|q|$.
\end{enumerate}
\end{lemma}

\begin{proof}
The zero-cost challenger always prepares, because $B_r(\mu)\ge g_L>0$.
Hence $e_x\ge\rho$.
The cost is independent of the class and of $R$.
Given the class $j$ and the order $x$, the challenger prepares with probability $e_x$, and then $\E[V_T]=t_j$.
Without preparation $\E[V_T]=t_0$.
This gives (a).
Average (a) over the class with weights $\mu_x$, $1-\mu_x$ to get (b).
Subtract (b) from (a) to get (c).
For (d), the investor's order is the observed order, so the price is $P(q)$ and the entry is $e_q$.
The payoff is $q[\E[V_T\mid j,q]-P(q)]-k|q|$.
Insert (c).
\end{proof}

\begin{lemma}[Price sufficiency]\label{l:suff}
Fix a candidate equilibrium.
If two orders $x\ne y$ have the same price $P(x)=P(y)$, then $\mu_x=\mu_y$.
The challenger's posterior at a price equals the dealer posterior $\mu_x$ of every order at that price.
\end{lemma}

\begin{proof}
The challenger sees only the price.
So the challenger uses the same entry probability $e\ge\rho>0$ after $x$ and after $y$.
Lemma~\ref{l:prices}(b) gives $P(x)-P(y)=e\,\Delta_T\,(\mu_x-\mu_y)$.
The left side is zero, $e>0$ and $\Delta_T>0$.
Hence $\mu_x=\mu_y$.
The challenger's posterior at the price is an average of the posteriors of the orders at that price.
All these posteriors are equal.
\end{proof}

Lemma~\ref{l:suff} lets the challenger use $B_r(\mu_x)$ at every price, also when prices pool orders, as in $\mathcal N$.
It covers arbitrary candidates.
The existence proofs below also check directly that the prices in $\mathcal U$ and $\mathcal F$ differ.

\subsection{The three trading outcomes}

Table~\ref{tab:tuples} states the three outcomes as explicit tuples.
The notation $\delta_q$ is the point mass at the order $q$.
The entry after an order is $e_x=\rho+(1-\rho)J_r\one\{x=1\}$ in the outcomes $\mathcal U$ and $\mathcal F$, and $e_x=\rho$ in $\mathcal N$.
Write $e_+=\rho+(1-\rho)J_r$.

\begin{table}[htbp]
\centering\small
\caption{The outcomes $\mathcal N$ (no informed trade), $\mathcal U$ (buys only), $\mathcal F$ (buys and sales). Orders $x=-1,0,1$.}\label{tab:tuples}
\begin{tabular}{@{}lccc@{}}
\toprule
 & $\mathcal N$ & $\mathcal U$ & $\mathcal F$\\
\midrule
$(\sigma_H,\sigma_L)$ & $(\delta_0,\delta_0)$ & $(\delta_1,\delta_0)$ & $(\delta_1,\delta_{-1})$\\
$s=\lambda\sigma_H(1)$, $t=\lambda\sigma_L(-1)$ & $(0,0)$ & $(\lambda,0)$ & $(\lambda,\lambda)$\\
$\Pr(X=x)$ & $(b,\,1-2b,\,b)$ & $(b,\,1-2b-\tfrac\lambda2,\,b+\tfrac\lambda2)$ & $(d,\,1-2d,\,d)$\\
$\mu_x$ & $(\tfrac12,\tfrac12,\tfrac12)$ & $(\tfrac12,\ \mu_0^{\mathcal U},\ M)$ & $(m,\tfrac12,M)$\\
$e_x$ & $(\rho,\rho,\rho)$ & $(\rho,\rho,e_+)$ & $(\rho,\rho,e_+)$\\
$P(x)$ & $t_0+\rho A_r(\tfrac12)$ (all $x$) & $t_0+\rho A_r(\tfrac12),\ t_0+\rho A_r(\mu_0^{\mathcal U}),\ t_0+e_+A_r(M)$ & $t_0+\rho A_r(m),\ t_0+\rho A_r(\tfrac12),\ t_0+e_+A_r(M)$\\
\bottomrule
\end{tabular}

\smallskip
\footnotesize In $\mathcal U$ the prices are listed in the order $x=-1,0,1$, and $\mu_0^{\mathcal U}=(1-2b-\lambda)/(2-4b-\lambda)<\tfrac12$.
\end{table}

In every outcome, bids are truthful.
A prepared challenger bids its value, and the incumbent bids $R$.
The challenger prepares at the price $P(x)$ if and only if it has cost $0$ or $B_r(\mu_x)\ge\cH$.
The prices in $\mathcal F$ and $\mathcal U$ are distinct across orders (Proposition~\ref{p:out}).
So the price reveals the order, and the entry rule is a function of $\mu_x$.

\begin{proposition}[Equilibrium outcomes; analytical]\label{p:out}
Let $\pi=1/2$ and $\beta\le1/3$.
\begin{enumerate}[label=(\alph*)]
\item If $\Delta_T(r)<k$, the outcome $\mathcal N$ is the unique equilibrium outcome. Entry equals $\rho$.
\item If $r_S<r<h$, the outcome $\mathcal F$ is the unique equilibrium outcome. Equivalently $\rho\Delta_T(r)\,m>k$.
\item If $r_F<r<r_N$ and $r\le r_C(M)$, the outcome $\mathcal U$ is an equilibrium outcome. The outcome $\mathcal N$ is also an equilibrium outcome for $r\le r_N$.
\item Part (b) holds without change if the challenger does not see the price. Then the challenger decides on the prior $\tfrac12$, entry equals $\rho$ after every order, and $J_r$ is replaced by $0$ in Table~\ref{tab:tuples}.
\end{enumerate}
\end{proposition}

\begin{proof}
\emph{Conditions.}
An equilibrium needs five conditions.
(i) Bayes consistency: $\mu_x$ is computed from the order law (Table~\ref{tab:tuples}).
(ii) Competitive prices: $P(x)=\E[V_T\mid x]$ (Lemma~\ref{l:prices}(b)).
(iii) Entry optimality: the challenger prepares if and only if its gross profit covers its cost.
(iv) Investor optimality: each class puts its mass on the orders with the largest $U_j(q)$ of Lemma~\ref{l:prices}(d).
(v) Truthful bids.
Conditions (i), (ii), (iii) and (v) hold by construction in each tuple.
Only (iv) needs a check.

(a) Take any candidate.
By Lemma~\ref{l:prices}(d), $U_H(1)\le\Delta_T-k<0$, $U_H(-1)<0$, and $U_H(0)=0$.
Also $U_L(-1)\le\Delta_T-k<0$, $U_L(1)<0$, and $U_L(0)=0$.
So both classes choose $0$.
All posteriors equal $\tfrac12$, and $B_r(\tfrac12)<\cH$ by Lemma~\ref{l:auction}(b).
Entry is $\rho$, and the tuple $\mathcal N$ satisfies (i) to (v).

(b) Take any candidate equilibrium.
Lemma~\ref{l:bounds} gives $\mu_1\le M$ and $\mu_{-1}\ge m$ for every profile.
Lemma~\ref{l:prices} gives $e_x\ge\rho$.
Hence
\[
U_H(1)\ge\rho\Delta_T(1-M)-k=\rho\Delta_T m-k>0,\qquad U_H(0)=0>U_H(-1),
\]
\[
U_L(-1)\ge\rho\Delta_T m-k>0,\qquad U_L(0)=0>U_L(1).
\]
The first inequality is the hypothesis $r>r_S$, because $\Delta_T(r)$ increases in $r$ and $\Delta_T(r_S)=k/(\rho m)$.
So $\sigma_H=\delta_1$ and $\sigma_L=\delta_{-1}$ in every equilibrium.
Bayes' rule then gives the posteriors $(m,\tfrac12,M)$.
The prices are $P(-1)=t_0+\rho A_r(m)<P(0)=t_0+\rho A_r(\tfrac12)<P(1)=t_0+e_+A_r(M)$, by Lemma~\ref{l:auction}(a) and $e_+\ge\rho$.
The prices are distinct, so the challenger infers $x$ and prepares as in Table~\ref{tab:tuples}.
This is an equilibrium, and it is the only one.

(c) Take $\mathcal U$.
The posteriors are $(\tfrac12,\mu_0^{\mathcal U},M)$ with $\mu_0^{\mathcal U}<\tfrac12$.
The condition $r\le r_C(M)$ gives $e_1=1$.
The orders $x=-1$ and $x=0$ do not recruit, because $B_r(\mu)\le B_r(\tfrac12)<\cH$.
Class $H$ earns $U_H(1)=\Delta_T(1-M)-k>0$ if and only if $r>r_F$.
It earns $U_H(0)=0$ and $U_H(-1)<0$.
Class $L$ earns $U_L(-1)=\rho\Delta_T/2-k<0$ if and only if $r<r_N$.
It earns $U_L(0)=0$ and $U_L(1)<0$.
The prices satisfy $P(0)<P(-1)<P(1)$, so they are distinct.
For $\mathcal N$, the deviation payoff of each class is $\rho\Delta_T/2-k$, which is not positive if and only if $r\le r_N$.

(d) With a hidden price, entry is $\rho$ after every order, because $B_r(\tfrac12)<\cH$.
The bounds $\mu_1\le M$, $\mu_{-1}\ge m$ and $e_x=\rho$ give the same inequalities as in (b).
\end{proof}

\begin{remark}[The endpoint $r_S$]\label{r:rS}
The bound in Proposition~\ref{p:out}(b) is the exact number $r_S$.
At $r=\reg{rS_counter_r_B}$, which lies in the interval $(\reg{rS_p5_B},h)$, the margin is $\rho\Delta_T(r)m-k=\reg{rS_counter_margin_B}$.
This margin is negative, so the low class strictly prefers to abstain in $\mathcal F$.
At $r=\reg{rS_cert_hi_B}$ the margin is positive, and it stays positive for larger $r$ because $\Delta_T$ increases.
The margin tends to zero as $r\downarrow r_S$, so no positive margin holds on the whole interval $(r_S,h)$.
At $r=r_S$ the low class is indifferent, and the outcome depends on the investor rule (Section~\ref{sec:margins}).
\end{remark}

\subsection{Observation protocol}\label{sec:protocol}

Two objects in the next sections are not part of the trading game.
They add observations and change no payoff.
\begin{itemize}
\item The \emph{pre-slot price} is $\bar P=\sum_x\Pr(X=x)P(x)$.
A risk-neutral dealer who sees no order sets this price, by iterated expectations.
This is the price before the slot under the passive observation protocol.
\item The \emph{preparation announcement} is a public event after the price $P(x)$.
It reveals whether the challenger prepared, and nothing else.
Dealers then reprice to $\E[V_T\mid X=x,\text{announcement}]$.
\end{itemize}
Prices are in units of the low value $\ell$.
The model sets the standalone value of the target to zero.
A return divides a price change by the price plus the standalone value.
So the ratios below are upper bounds in size for any positive standalone value.
Both dollar changes and ratios to $P(x)$ appear in the tables.

Table~\ref{tab:objects} links each model object to its empirical counterpart and states the condition for the link.

\begin{table}[H]
\centering\small
\caption{Model objects and their empirical counterparts.}\label{tab:objects}
\begin{tabular}{@{}p{3.0cm}p{9.9cm}@{}}
\toprule
Model object & Empirical counterpart and condition\\
\midrule
$P(x)=\E[V_T\mid X=x]$ & The conditional execution price of a unit order. The model has no quote stage and no limit order book. Zero dealer profit gives the price formula. A precommitted bid and ask need an added institutional assumption.\\
$P(1)-P(-1)$ & The difference between the conditional execution prices of a unit buy and a unit sale.\\
$X=0$ & A public opportunity with no trade. Keep the opportunity as the denominator.\\
$v=\Pr(X\ne0)$ & The number of nonzero unit trades per public opportunity. It is not volume per unit of calendar time.\\
$P(0)$ & The conditional value after an observed opportunity with no trade. It is not the price before the slot.\\
$I_-=P(0)-P(-1)$ & The sale price difference relative to the price after no trade. The measurement needs this reference price.\\
$\bar P=\E[P(X)]$ & The price before the slot, under the passive observation protocol.\\
$A$ & Preparation before the auction. An observed final bid is a different selection event.\\
$w_x$ & The win probability of the challenger given preparation and the order. Keep it apart from $e_xw_x$.\\
$\Delta^a_x$ & The announcement price change, when the event reveals only preparation, on events of positive probability.\\
$\Delta^a_x/P(x)$ & The simple return when the standalone value is zero. A positive standalone value enlarges the denominator.\\
$\vartheta$ & The share of nonzero orders that come from the informed identity on $\mathcal F$. It differs from the presence probability $\lambda$.\\
\bottomrule
\end{tabular}
\end{table}

The first measurement rule is this.
Use the first prespecified public opportunity and record a zero when no trade occurs.
A sample of first executed trades omits $X=0$.
Such a sample cannot identify $v$, the price effect of a neutral order, or the informed share from the formulas below.

The friction $k$ enters the investor's choice.
It does not enter the execution spread as a dealer charge.
The sell-impact formula of Proposition~\ref{p:ident}(e) uses maintained inputs.
These inputs are the entry floor $\rho$, or its estimate from non-buys, and the value $\Delta_T(r)$ from the public class values and the incumbent strength.
The formula identifies the informed share inside the maintained auction model.
It is not a model-free estimator.

\subsection{Calibration B and the arithmetic}

Calibration B sets $1+u=e$ and $\beta=0.294655$, an exact decimal.
Then $\lambda=u\beta/(1+u\beta)=\reg{lambda_p5_B}$, $b=\reg{b_p5_B}$, $d=\reg{d_slot_B}$, $M=\reg{M_p5_B}$, $m=\reg{m_p5_B}$, and $\vartheta=\reg{theta_B}$.
These match the starter values of the brief.
The script lifts $e$ to the enclosure \texttt{iv.e} (60 digits) and $\beta$ to the interval $294655/10^6$.
Every sign test in the proofs is an interval test.
A test passes only if the enclosure excludes zero.
The run log prints the enclosures.
The script has no assertion and no test.
It reports an undecided sign as a flag; the log shows none.

\section{Scope of the results}\label{sec:scope}

The package proves its results for the outcomes $\mathcal N$, $\mathcal U$ and $\mathcal F$ at the prior $1/2$.
These outcomes cover the three benchmark strengths and the buy-only strengths $1.55$, $1.60$ and $1.65$.
Three items stay open.
Table~\ref{tab:scope} states, for each item, what the package proves and what is open.
This table is the only place that lists the open items.

\begin{table}[H]
\centering\small
\caption{Scope: what the package proves and what stays open.}\label{tab:scope}
\begin{tabular}{@{}p{2.5cm}p{7.2cm}p{3.2cm}@{}}
\toprule
Item & What the package proves & Open\\
\midrule
Equilibrium selection & For $r_F<r<r_N$ and $r\le r_C(M)$, the equilibrium outcomes are exactly $\mathcal N$ and $\mathcal U$ under the challenger rule (Proposition~\ref{p:coex}(v)). Proposition~\ref{p:coex} states the observables of both. An indifferent challenger who may mix adds one partial-buy outcome (Remark~\ref{r:mix}). & Which outcome occurs.\\
Intermediate branch & For $\Delta_T(r)<k$ the unique outcome is $\mathcal N$. For $r>r_S$ the unique outcome is $\mathcal F$. On $[r_N,r_S]$ the low class may sell at a rate that makes it indifferent. Proposition~\ref{p:entry}(iv) gives the entry rate for every low-class rate, and Lemma~\ref{l:law} gives the order law. & The equilibrium set for $\Delta_T(r)\ge k$ with $r\le r_F$, and the spread, winner and announcement formulas on $[r_N,r_S]$.\\
Observation protocols & The formulas hold under the protocols of Section~\ref{sec:protocol}: the opportunity denominator, conditional execution prices, the passive pre-slot price, announcements that reveal only preparation, and a zero standalone value. & A data source for each protocol.\\
\bottomrule
\end{tabular}
\end{table}

\section{Order-flow observables}\label{sec:flow}

Let $s=\lambda\Pr(q=1\mid H,\text{informed})$ and $t=\lambda\Pr(q=-1\mid L,\text{informed})$.
By Lemma~\ref{l:prices}(d), a high-class sale and a low-class purchase lose money at every schedule.
So no equilibrium contains them.
Let $v=\Pr(X\ne0)$ be the volume.

\begin{lemma}[Order law]\label{l:law}
At $\pi=\tfrac12$, for $s,t\in[0,\lambda]$:
\[
\Pr(X=1)=b+\tfrac s2,\quad \Pr(X=-1)=b+\tfrac t2,\quad v=2b+\tfrac{s+t}2,\quad \E[X]=\tfrac{s-t}{2}.
\]
The informed share of nonzero orders is $(s+t)/(4b+s+t)$.
The informed share among buys is $s/(2b+s)=2\mu_1-1$.
The informed share among sales is $t/(2b+t)=1-2\mu_{-1}$.
\end{lemma}

\begin{proof}
Class $H$ places a buy with mass $b+s$ and class $L$ with mass $b$.
Class $H$ places a sale with mass $b$ and class $L$ with mass $b+t$.
Average over the classes to get the order law.
An informed trade has probability $(s+t)/2$ per slot.
Divide by $v$ for the share.
For buys, the informed mass is $s/2$ and the buy mass is $b+s/2$; the ratio is $s/(2b+s)$.
The posterior is $\mu_1=(b+s)/(2b+s)$, so $2\mu_1-1=s/(2b+s)$.
The sale case is the same.
\end{proof}

\begin{proposition}[Order-flow signature; analytical, values computer-assisted]\label{p:flow}
The outcomes of Table~\ref{tab:tuples} have the following order-flow observables.
\begin{center}\small
\begin{tabular}{@{}lccc@{}}
\toprule
& $\mathcal N$ & $\mathcal U$ & $\mathcal F$\\
\midrule
volume $v$ & $2b$ & $2b+\lambda/2$ & $2b+\lambda$\\
mean order $\E[X]$ & $0$ & $\lambda/2$ & $0$\\
informed share of nonzero orders & $0$ & $u/(4+u)$ & $u/(2+u)=2M-1$\\
\bottomrule
\end{tabular}
\end{center}
Consequently:
\begin{enumerate}[label=(\roman*)]
\item $v_{\mathcal N}<v_{\mathcal U}<v_{\mathcal F}$.
\item $\E[X]>0$ in $\mathcal U$ and $\E[X]=0$ in $\mathcal N$ and in $\mathcal F$.
\item The informed share in $\mathcal U$ is below the share in $\mathcal F$.
In $\mathcal U$ it equals the normalized imbalance $\E[X]/v$.
\item On $r_S<r<10$, the order law, the volume, the mean order, the informed share and the posterior vector $(m,\tfrac12,M)$ do not depend on $r$.
\end{enumerate}
\end{proposition}

\begin{proof}
Insert $(s,t)=(0,0),(\lambda,0),(\lambda,\lambda)$ into Lemma~\ref{l:law}.
For $\mathcal U$: $(s+t)/(4b+s+t)=\lambda/(4b+\lambda)=u/(4+u)$, because $u=\lambda/b$.
This is below $u/(2+u)$ because $4+u>2+u$.
Also $\E[X]/v=(\lambda/2)/(2b+\lambda/2)=\lambda/(4b+\lambda)$.
Statement (i) follows from $\lambda>0$ and $b>0$.
Statement (iv) holds because the tuple $\mathcal F$ of Proposition~\ref{p:out}(b) contains no $r$ in its order law or its beliefs.
\end{proof}

\paragraph{At B.}
The volumes are \reg{vol_none_B}, \reg{vol_buy_B}, \reg{vol_full_B}.
The mean order in $\mathcal U$ is \reg{imb_buy_B}.
The informed shares are \reg{theta_buy_B} in $\mathcal U$ and \reg{theta_B} in $\mathcal F$.
The earlier study reported $9/31$ and $9/20$ at $\beta=1/3$.

\paragraph{Observable.}
Use the first prespecified public opportunity after the sale becomes public, and record a zero when no trade occurs.
Volume is the number of nonzero orders per opportunity.
The informed share is a parameter of the order law and not an estimate of an external PIN measure.

\section{Prices and the spread}\label{sec:spread}

Let $S=P(1)-P(-1)$ be the conditional execution spread, $I_+=P(1)-P(0)$ the buy impact, and $I_-=P(0)-P(-1)$ the sell impact.
The price $P(x)$ is the expected target value given the order $x$.
It is also the price at which a dealer earns zero profit on that order.
So $S$ is the difference between the conditional execution prices of a unit buy and a unit sale.
The model has no separate quote stage.

\begin{proposition}[Spread decomposition; analytical, values computer-assisted]\label{p:spread}
Let $r_S<r<h$ (outcome $\mathcal F$) and $M>5/8$.
\begin{enumerate}[label=(\roman*)]
\item $P(-1)=t_0+\rho A_r(m)$, $P(0)=t_0+\rho A_r(\tfrac12)$, $P(1)=t_0+\rho A_r(M)+(1-\rho)J_rA_r(M)$.
\item $S=\underbrace{\rho\Delta_T\vartheta}_{\text{class component}}+\underbrace{(1-\rho)J_rA_r(M)}_{\text{entry component}}$.
\item $I_-=\rho\Delta_T\vartheta/2$ and $I_+=I_-+(1-\rho)J_rA_r(M)$. Hence $I_+>I_-$ if and only if $J_r=1$.
\item If the challenger does not see the price, the unique outcome is $\mathcal F$ and the spread equals the class component $\rho\Delta_T\vartheta$. The ratio of the two spreads is $1+(1-\rho)J_rA_r(M)/(\rho\Delta_T\vartheta)$.
\item Let $r_S<r_C(M)$. On $(r_S,r_C]$, the class component, the entry component, $I_-$, $I_+$ and $S$ all increase strictly in $r$.
On $(r_C,h)$, the entry component is identically zero. The class component, $I_-$, $I_+=I_-$ and $S$ increase strictly there.
\item Let $r_S<r_C(M)$, so that the ceiling lies inside the interval of $\mathcal F$ (Lemma~\ref{l:auction}(d) gives $r_C<h$). At $r_C$ (challenger tie rule: prepare), $S(r_C)=\rho\Delta_T\vartheta+(1-\rho)A_{r_C}(M)$. The right limit is $S(r_C^+)=\rho\Delta_T\vartheta$. The spread falls by $(1-\rho)A_{r_C}(M)>0$. The price $P(0)$, the price $P(-1)$ and $I_-$ are continuous at $r_C$. The price $P(1)$ and $I_+$ fall by the same amount. The pre-slot price $\bar P$ falls by $(1-\rho)d\,A_{r_C}(M)$.
\end{enumerate}
\end{proposition}

\begin{proof}
(i) is Table~\ref{tab:tuples} with Lemma~\ref{l:prices}(b).
(ii) Subtract $P(-1)$ from $P(1)$.
The term $t_0$ cancels.
The remaining terms are $\rho[A_r(M)-A_r(m)]+(1-\rho)J_rA_r(M)$.
Since $A_r$ has slope $\Delta_T$, $A_r(M)-A_r(m)=\Delta_T(M-m)=\Delta_T\vartheta$.
(iii) Likewise $A_r(\tfrac12)-A_r(m)=\Delta_T(\tfrac12-m)=\Delta_T\vartheta/2$ and $A_r(M)-A_r(\tfrac12)=\Delta_T\vartheta/2$.
The buy impact adds the entry term.
(iv) Proposition~\ref{p:out}(d) gives $\mathcal F$ with entry $\rho$.
Insert $J_r=0$ in (ii).
(v) Within an interval where $J_r$ is constant, $\Delta_T$ and $A_r(M)$ increase strictly by Lemma~\ref{l:auction}(a).
The factors $\rho$, $\vartheta$, $1-\rho$ do not depend on $r$.
On $(r_S,r_C]$, $J_r=1$, so the entry component $(1-\rho)A_r(M)$ increases strictly, and so do $I_-$, $I_+$ and $S$ by (ii) and (iii).
On $(r_C,h)$, $J_r=0$ by Lemma~\ref{l:auction}(d).
The entry component is then $0$ at every $r$, and its derivative is zero.
By (ii) and (iii), $I_-$, $I_+=I_-$ and $S=\rho\Delta_T\vartheta$ increase strictly there, because $\Delta_T$ does.
(vi) By Lemma~\ref{l:auction}(d), $J_r=1$ at $r=r_C$ and $J_r=0$ for $r>r_C$.
All other terms are continuous in $r$.
The price $P(1)$ loses $(1-\rho)A_{r_C}(M)$.
In $\bar P$, the order $x=1$ has probability $d$ and the other terms do not change.
\end{proof}

\paragraph{At B.}
Table~\ref{tab:spread} gives the components.
The entry component is \reg{ba_share_r1_B} of the spread at $r_1$.
With a hidden price the spread at $r_1$ would be \reg{ba_class_r1_B}.
Price access multiplies the spread by \reg{ba_ratio_r1_B}.
The ceiling is $r_C=\reg{rC_check_B}$, which lies \reg{margin_obs_r2_above_rC_B} below $r_2$ and \reg{margin_obs_rS_rC_B} above $r_S$.
At $r_C$ the spread falls by \reg{ba_jump_rC_B}, from \reg{ba_rC_left_B} to \reg{ba_rC_right_B}.
The expected price $\bar P$ falls by \reg{pbar_jump_rC_B} (change of sign included).

\begin{table}[htbp]
\centering\small
\caption{Spread components at B (units of $\ell$).}\label{tab:spread}
\begin{tabular}{@{}lcc@{}}
\toprule
 & $r_1=3$ & $r_2=18/5$\\
\midrule
class component $\rho\Delta_T\vartheta$ & \reg{ba_class_r1_B} & \reg{ba_r2_B}\\
entry component $(1-\rho)J_rA_r(M)$ & \reg{ba_entry_r1_B} & $0$\\
spread $S$ & \reg{ba_r1_B} & \reg{ba_r2_B}\\
sell impact $I_-$ & \reg{impact_sell_r1_B} & \reg{impact_sell_r2_B}\\
buy impact $I_+$ & \reg{impact_buy_r1_B} & \reg{impact_sell_r2_B}\\
$I_+-I_-$ & \reg{impact_asym_r1_B} & $0$\\
\bottomrule
\end{tabular}
\end{table}

\paragraph{Observable.}
Take the execution price of a unit buy minus the execution price of a unit sale, in the target's shares.
Take the buy and sell price impacts relative to the price after a recorded zero.
The sign of $I_+-I_-$ marks the processes in which a buy moves the entry of the challenger.
This test needs no knowledge of the auction objects.
Above $r_C$ the entry component is zero, so $I_+-I_-=0$ there.

\section{Identification of the informed share}\label{sec:ident}

\begin{proposition}[Identification; analytical, values computer-assisted]\label{p:ident}
Let $r>r_S$ (outcome $\mathcal F$), and write $v=2b+\lambda$.
\begin{enumerate}[label=(\alph*)]
\item \emph{Counts.} The order law is $(v/2,1-v,v/2)$. It depends on $(\lambda,\beta)$ only through $v$.
For each $u'>0$ there is exactly one pair $(\lambda',\beta')$ with odds $u'$ and the same $v$:
\[
\lambda'=\frac{vu'}{u'+2},\qquad b'=\frac{v}{u'+2},\qquad \beta'=\frac{v}{2+u'(1-v)}.
\]
The condition $\beta'\le 1/3$ holds if and only if $u'\ge u_{\min}(v)=\max\{0,(3v-2)/(1-v)\}$.
The informed share of this pair is $\vartheta'=u'/(u'+2)$.
This part describes order laws, before investor optimality.
\item \emph{Entry.} Entry in $\mathcal F$ equals $\rho+(1-\rho)J_rv/2$. It depends on $(\lambda,\beta)$ only through $v$ and $J_r$.
\item \emph{Incentives.} Let $\bar\vartheta(r)=1-2k/(\rho\,\Delta_T(r))$.
A full-order pair with share $\vartheta'$ gives its counts through an $\mathcal F$ equilibrium at the strength $r$ only if $\vartheta'\le\bar\vartheta(r)$.
The pair has $\mathcal F$ as its unique outcome if $\vartheta'<\bar\vartheta(r)$.
The bound $\bar\vartheta(r)$ increases in $r$.
Let $\vartheta_{\min}(v)=u_{\min}(v)/(u_{\min}(v)+2)$ and assume $\vartheta_{\min}(v)<\bar\vartheta(r)$.
Then the shares that give the counts of $v$ through an $\mathcal F$ equilibrium at $r$ form the interval $[\vartheta_{\min}(v),\bar\vartheta(r))$.
At the endpoint $\bar\vartheta(r)$ the outcome depends on the investor rule.
\item \emph{Band.} Suppose $B_{r_1}(M')>\cH>B_{r_2}(M')$ for the pair $(\lambda',\beta')$.
Then $\vartheta'\in(2\tau(r_1)-1,\,2\tau(r_2)-1)$.
If $v\le200/259$, every $\vartheta'$ in this interval comes from a pair with the same $v$.
These pairs have the same entry at $r_1$ and the same entry at $r_2$.
Each pair has $\mathcal F$ as its unique outcome at $r_1$ and at $r_2$, with the margin $\rho\Delta_T(r_i)(1-\vartheta')/2-k>\rho\Delta_T(r_1)(1-\tau(r_2))-k>0$.
\item \emph{Prices.} The entry floor $\rho$ equals $\Pr(\text{prepare}\mid X\ne1)$. The sell impact identifies the informed share and the slot parameters:
\[
\vartheta=\frac{2I_-}{\rho\,\Delta_T(r)},\qquad \lambda=v\vartheta,\qquad
\beta=\frac{v(1-\vartheta)}{2(1-v\vartheta)},\qquad u=\frac{2\vartheta}{1-\vartheta}.
\]
\item \emph{Buy-only branch.} In $\mathcal U$ the counts identify the slot: $\lambda=2\E[X]$, $b=(v-\E[X])/2$, $\vartheta_{\mathcal U}=\E[X]/v$ and $\beta=b/(1-\lambda)$.
\end{enumerate}
\end{proposition}

\begin{proof}
(a) Lemma~\ref{l:law} with $s=t=\lambda$ gives $\Pr(X=\pm1)=d=v/2$.
Given $v$ and $u'$, the identities $v=\lambda'+2b'$ and $u'=\lambda'/b'$ give $b'=v/(u'+2)$ and $\lambda'=u'b'$.
Then $\beta'=b'/(1-\lambda')=v/(u'+2-vu')$.
The inequality $\beta'\le1/3$ is $3v\le 2+u'(1-v)$.
Solve for $u'$, using $v<1$.

(b) Entry is $\sum_x\Pr(X=x)e_x=\rho+(1-\rho)J_r\Pr(X=1)$.

(c) In $\mathcal F$ the sale has the posterior $m'=(1-\vartheta')/2$.
A sale does not recruit, because $m'\le\tfrac12$ and $B_r(\mu)\le B_r(\tfrac12)<\cH$ for $\mu\le\tfrac12$.
So the entry after a sale is $\rho$, and Lemma~\ref{l:prices}(d) gives $U_L(-1)=\rho\Delta_T m'-k$.
This is nonnegative if and only if $\vartheta'\le\bar\vartheta(r)$.
If it is negative, the low class strictly prefers to abstain, and $\mathcal F$ is not an equilibrium.
Now let $\vartheta'<\bar\vartheta(r)$.
The posterior bounds of the pair (Lemma~\ref{l:bounds} with the odds $u'$) give $\mu_1\le M'$ and $\mu_{-1}\ge m'$ for every profile.
The entry is at least $\rho$ at every order.
As in the proof of Proposition~\ref{p:out}(b), $U_H(1)\ge\rho\Delta_T m'-k>0$ and $U_L(-1)\ge\rho\Delta_T m'-k>0$.
The wrong-sign orders lose, so $\mathcal F$ is the unique outcome of the pair.
The bound $\bar\vartheta(r)$ increases in $r$, because $\Delta_T$ does.
The interval for the shares combines this bound with part (a): $\beta'\le1/3$ holds if and only if $u'\ge u_{\min}(v)$.

(d) The pair has the same $v$, so the same order law, and by (b) the same entry given $J_r$.
The condition $J_{r_1}=1$ is $M'\ge\tau(r_1)$ and $J_{r_2}=0$ is $M'<\tau(r_2)$.
Since $\vartheta'=2M'-1$, the first claim follows.
For the second claim, $u_{\min}(v)\le 82/59$ if and only if $v\le 200/259$, and $82/59$ is the lower end of the band for $u'$ (the value $\tau(r_1)=141/200$ gives $u=82/59$).
For the third claim, $\vartheta'<2\tau(r_2)-1$ gives $m'=(1-\vartheta')/2>1-\tau(r_2)$.
Hence $\rho\Delta_T(r_1)m'-k>\rho\Delta_T(r_1)(1-\tau(r_2))-k=21563/870600$, and this number is positive.
The value $\Delta_T(r_2)>\Delta_T(r_1)$ gives the same bound at $r_2$.
Part (c) then gives the unique outcome $\mathcal F$ at both strengths.

(e) Entry after $x\ne1$ equals $\rho$ in $\mathcal F$.
Proposition~\ref{p:spread}(iii) gives $I_-=\rho\Delta_T\vartheta/2$.
Solve for $\vartheta$.
Then $v=\lambda+2b$ and $\vartheta=\lambda/v$ give $\lambda=v\vartheta$ and $b=v(1-\vartheta)/2$.
Finally $\beta=b/(1-\lambda)$ and $u=2\vartheta/(1-\vartheta)$ from $\vartheta=u/(2+u)$.

(f) In $\mathcal U$, $\Pr(X=1)-\Pr(X=-1)=\lambda/2$ and $v=2b+\lambda/2$.
Solve for $\lambda$ and $b$.
Then $\vartheta_{\mathcal U}=u/(4+u)=\lambda/(4b+\lambda)=\E[X]/v$.
\end{proof}

\paragraph{At B.}
The volume is $v=\reg{vol_full_B}$.
The pair with the smallest admissible odds ($\beta'=1/3$) has informed share \reg{theta_min_counts_B}.
Order laws alone allow every share in $[\reg{theta_min_counts_B},1)$.
The investor incentives cut this range.
An $\mathcal F$ equilibrium needs $\vartheta'<\reg{theta_max_full_eq_r1_B}$ at $r_1$ and $\vartheta'<\reg{theta_max_full_eq_r2_B}$ at $r_2$.
So the counts allow the shares in $[\reg{theta_min_counts_B},\reg{theta_max_full_eq_r1_B})$ through an $\mathcal F$ equilibrium at both strengths.
The share $\vartheta'=0.9$ has the odds $u'=18$ and the pair $\beta'=\reg{theta09_beta_B}$, with the same counts as B.
The low class would earn \reg{theta09_lowsale_r1_B} from a full sale at $r_1$.
So this pair gives these counts in no $\mathcal F$ equilibrium at $r_1$.
Inside the band defined by $B_{r_1}(M')>\cH>B_{r_2}(M')$, the share lies in $(\reg{theta_lo},\reg{theta_hi})$.
The incentive margin on the whole band is \reg{margin_obs_band_incentive_B}.
The pair $u'=8/5$ has $\lambda'=\reg{alt_lambda_B}$, $\beta'=\reg{alt_beta_B}$, and $\vartheta'=\reg{alt_theta_B}$.
It has the volume \reg{vol_full_B} and the entry \reg{alt_entry_r1_B} at $r_1$, the same as B.
Its spread at $r_1$ is \reg{alt_ba_r1_B}, against \reg{ba_r1_B} at B.
So prices separate the two pairs.
The identification formulas of (e) recover $(\vartheta,\lambda,\beta)=(\reg{theta_B},\reg{lambda_p5_B},\reg{beta_p5_B})$ from the model sell impact \reg{impact_sell_r1_B}, with differences below $10^{-50}$ in the run log.

\paragraph{Observable.}
The inputs are the order frequencies, including recorded zeros, and the sell impact relative to the price after no trade.
The other inputs are the entry rate after non-buys and $\Delta_T(r)$.
The value $\Delta_T(r)$ comes from the public strength and the class values.
The entry floor $\rho$ is a maintained input, or an estimate from the non-buys.
The informed share is a parameter of this model.

\section{Entry across strengths}\label{sec:entry}

Let $A$ be the event that the challenger prepares.
Let $\Ent=\Pr(A)$.

\begin{proposition}[Entry; analytical, values computer-assisted]\label{p:entry}
Let $\pi=\tfrac12$ and $M>5/8$.
\begin{enumerate}[label=(\roman*)]
\item \emph{Entry by order.} In $\mathcal F$ and $\mathcal U$, $\Pr(A\mid X=x)=e_x$ with $e_x=\rho$ for $x\ne1$ and $e_1=\rho+(1-\rho)J_r$. In $\mathcal N$, $e_x=\rho$.
\item \emph{Entry and buys.} In $\mathcal F$ and $\mathcal U$, $\Ent=\rho+(1-\rho)J_r\Pr(X=1)$, and $\Pr(A\mid X=1)-\Pr(A\mid X\ne1)=(1-\rho)J_r$.
\item \emph{Plateau.} Both $\mathcal F$ and $\mathcal U$ give $\Pr(X=1)=d$. So $\Ent=\rho+(1-\rho)d$ for $r_S<r\le r_C$ in $\mathcal F$, and for $r_F<r<r_N$ with $r\le r_C$ in $\mathcal U$. In $\mathcal F$, $\Ent=\rho$ for $r_C<r<h$. In $\mathcal N$, $\Ent=\rho$.
\item \emph{Any low-class behavior.} Let the high class buy with probability one and let $J_r=1$. Then $\Ent=\rho+(1-\rho)d$ for every low-class sale rate $t\in[0,\lambda]$.
\item \emph{Hidden price.} If the challenger does not see the price, $\Ent=\rho$ at every strength in $(\ell,h)$ and for every order profile.
\item \emph{Comparative statics.} On a region of $(\lambda,\beta)$ where $J_r=1$, $d=\beta+\lambda(\tfrac12-\beta)$ strictly increases in $\lambda$ and in $\beta$.
\end{enumerate}
\end{proposition}

\begin{proof}
(i) and (ii) follow from Table~\ref{tab:tuples}.
(iii) In $\mathcal F$, $\Pr(X=1)=b+\lambda/2$.
In $\mathcal U$, $\Pr(X=1)=b+\lambda/2$ as well.
The ceiling is Lemma~\ref{l:auction}(d).
(iv) A sale rate $t$ changes $\Pr(X=-1)$ and $\mu_{-1}=b/(2b+t)\le\tfrac12$ and $\mu_0\le\tfrac12$.
The buy posterior $\mu_1=(b+\lambda)/(2b+\lambda)=M$ does not contain $t$, and $B_r(\mu_{-1}),B_r(\mu_0)\le B_r(\tfrac12)<\cH$.
So only the order $x=1$ recruits, with probability $b+\lambda/2$.
(v) The challenger's posterior is the prior $\tfrac12$ after every order, whatever the order profile.
Lemma~\ref{l:auction}(b) gives $B_r(\tfrac12)<\cH$ for every $r\in(\ell,h)$.
Only the challenger with cost $0$ prepares, so $e_x=\rho$ for every order.
(vi) $\partial d/\partial\lambda=\tfrac12-\beta>0$ and $\partial d/\partial\beta=1-\lambda>0$.
\end{proof}

\paragraph{At B.}
The entry is $\rho=\reg{entry_floor_obs}$ at $r_0$.
It is \reg{entry_plateau_B} at $r_1$, a rise of \reg{entry_jump_B}.
It is $\reg{entry_r2_obs_B}$ at $r_2$.
The plateau \reg{U_entry_B} also holds on the buy-only branch at 1.55, 1.60 and 1.65.
With a hidden price, entry is \reg{entry_floor_obs} at $r_0$, $r_1$, $r_2$.

\begin{proposition}[Predictions of Section~7 of the Laplace manuscript in the slot; analytical, values computer-assisted]\label{p:sec7}
Let the brief's hypotheses hold: $\Delta_T(r_0)<k$, $\rho\Delta_T(r_1)m>k$, $B_{r_1}(M)>\cH>B_{r_2}(M)$.
Then $\rho\Delta_T(r_2)m>\rho\Delta_T(r_1)m>k$ also holds, because $\Delta_T$ increases in $r$.
\begin{enumerate}[label=(\arabic*)]
\item \emph{Strength raises entry when the price is visible.} $\Ent(r_1)-\Ent(r_0)=(1-\rho)d>0$. With a hidden price, $\Ent(r_1)=\Ent(r_0)=\rho$.
\item \emph{Favorable orders predict winners among prepared challengers.} See Proposition~\ref{p:win}.
\item \emph{Strong competition stops the entry response but not the information.} $\Ent(r_2)=\rho$ and $\Pr(A\mid X=x)=\rho$ for all $x$. The sell impact stays positive at $r_2$. The order law and the posterior vector equal those at $r_1$.
\end{enumerate}
\end{proposition}

\begin{proof}
(1) Proposition~\ref{p:out}(a) gives $\mathcal N$ at $r_0$ with entry $\rho$.
Proposition~\ref{p:out}(b) gives $\mathcal F$ at $r_1$ and at $r_2$, because $r_2>r_1>r_S$.
Proposition~\ref{p:entry}(iii) gives $\Ent(r_1)=\rho+(1-\rho)d$ since $J_{r_1}=1$.
The hidden-price case is Proposition~\ref{p:entry}(v).
(3) $J_{r_2}=0$, so Proposition~\ref{p:entry}(i) gives entry $\rho$ at every order.
Proposition~\ref{p:flow}(iv) gives the unchanged order law and posteriors.
Proposition~\ref{p:spread}(iii) gives $I_-=\rho\Delta_T\vartheta/2>0$.
\end{proof}

\paragraph{Observable.}
The empirical counterpart of preparation is the start of substantive diligence or a proposal that needs preparation.
Incumbent strength is measured from information public before the decision.
The order is the outcome of the first prespecified public opportunity after the sale becomes public, dated before the decision.
A recorded zero counts as the order $X=0$.

\section{Winners}\label{sec:win}

Let $W$ be the event that the challenger wins the target.
Let $w_x=\Pr(W\mid X=x,A)$.

\begin{proposition}[Winners; analytical, values computer-assisted]\label{p:win}
In the point-mass benchmark ($h>r>\ell>p$):
\begin{enumerate}[label=(\roman*)]
\item $w_x=\mu_x+(1-\mu_x)\ell/r$, and $\Pr(W\mid X=x)=e_xw_x$.
\item In each outcome of Table~\ref{tab:tuples}, the order with the higher price has the higher posterior and the higher $w_x$. Among prepared challengers, the win probability rises strictly with the price in $\mathcal U$ and $\mathcal F$.
\item In $\mathcal F$, $w_1-w_{-1}=\vartheta(1-\ell/r)$ and $w_1-w_0=w_0-w_{-1}=\vartheta(1-\ell/r)/2$. The gap rises strictly with $r$.
\item In $\mathcal F$, $\E[X]=0$, so $\Cov(X,\one_A)=(1-\rho)J_rd$.
For $r\le r_C$, $\Cov(X,\one_W)=d[w_1-\rho w_{-1}]$.
For $r>r_C$, $\Cov(X,\one_A)=0$ and
\[
\Cov(X,\one_W)=d\,\rho\,\vartheta\,(1-\ell/r)>0 .
\]
\item The same covariances hold at every $r>r_S$ with a hidden price: $\Cov(X,\one_A)=0$ and $\Cov(X,\one_W)=d\rho\vartheta(1-\ell/r)$.
\end{enumerate}
\end{proposition}

\begin{proof}
(i) The cost is independent of the class and of $R$.
So preparation does not change the class posterior: $\Pr(H\mid X=x,A)=\mu_x$.
A prepared class-$H$ challenger wins, because $h>r$.
A prepared class-$L$ challenger wins if and only if $R<\ell$, with probability $\ell/r$.
This gives $w_x$.
Multiply by $e_x$ for the unconditional probability.
(ii) In $\mathcal F$ and $\mathcal U$ the prices are strictly ordered like the posteriors (proof of Proposition~\ref{p:out}).
The function $\mu\mapsto\mu+(1-\mu)\ell/r$ is strictly increasing because $\ell<r$.
(iii) Subtract: $w_1-w_{-1}=(\mu_1-\mu_{-1})(1-\ell/r)=\vartheta(1-\ell/r)$.
The derivative in $r$ is $\vartheta\ell/r^2>0$.
(iv) Because $\E[X]=0$, $\Cov(X,Y)=\E[XY]=d(y_1-y_{-1})$ with $y_x=\E[Y\mid X=x]$.
For $Y=\one_A$, $y_1-y_{-1}=e_1-\rho=(1-\rho)J_r$.
For $Y=\one_W$, $y_x=e_xw_x$.
If $r>r_C$, $e_x=\rho$, so $y_1-y_{-1}=\rho(w_1-w_{-1})$.
(v) Replace $e_x$ by $\rho$ in (iv).
\end{proof}

\paragraph{At B.}
The win probability among prepared challengers after a sale, a neutral order and a buy is (\reg{wprep_sell_r1_B}, \reg{wprep_neutral_r1_B}, \reg{wprep_buy_r1_B}) at $r_1$.
It is (\reg{wprep_sell_r2_B}, \reg{wprep_neutral_r2_B}, \reg{wprep_buy_r2_B}) at $r_2$.
The winner gap is \reg{wgap_r1_B} at $r_1$ and \reg{wgap_r2_B} at $r_2$.
The covariance of the order with entry is \reg{cov_entry_r1_B} at $r_1$ and \reg{cov_entry_r2_B} at $r_2$.
The covariance of the order with a challenger win is \reg{cov_win_r1_B} at $r_1$ and \reg{cov_win_r2_B} at $r_2$.
The challenger wins with probability \reg{win_prob_r1_B} at $r_1$ and \reg{win_prob_r2_B} at $r_2$.

\paragraph{Observable.}
Take processes with a prepared further buyer.
Regress the identity of the winner on the signed order or on the price.
The slope is positive.
Above $r_C$ the same regression on entry has slope zero.
The gap $w_1-w_{-1}$ concerns prepared challengers.
The unconditional gap $\Pr(W\mid X=1)-\Pr(W\mid X=-1)=e_1w_1-\rho w_{-1}$ also contains entry.
Above $r_C$ it equals $\rho(w_1-w_{-1})$.

\section{Event returns}\label{sec:event}

Use the observation protocol of Section~\ref{sec:setup}.
Write $\Delta^1_x$ for the price change when the announcement says the challenger prepared, and $\Delta^0_x$ for the change when it says the challenger did not prepare.

\begin{proposition}[Announcement returns; analytical, values computer-assisted]\label{p:event}
In every outcome of Table~\ref{tab:tuples}:
\begin{enumerate}[label=(\roman*)]
\item $\Delta^1_x=(1-e_x)A_r(\mu_x)$. If $e_x<1$, then $\Delta^0_x=-e_xA_r(\mu_x)$.
\item If $e_x<1$, then $e_x\Delta^1_x+(1-e_x)\Delta^0_x=0$. The mean is taken over both announcements, with their model probabilities. The mean of the ratios $\Delta^a_x/P(x)$ is also zero, because $P(x)$ is fixed given $x$.
\item $\Delta^1_x=0$ if and only if $e_x=1$. Also $\Pr(A^c\mid X=x)=0$ if and only if $e_x=1$. At $e_x=1$ no posterior after $A^c$ exists, so $\Delta^0_x$ has no observed counterpart.
\item In $\mathcal F$ with $J_r=1$, the preparation changes after a sale, a neutral order and a buy are $((1-\rho)A_r(m),\,(1-\rho)A_r(\tfrac12),\,0)$. In $\mathcal F$ with $J_r=0$ they are $((1-\rho)A_r(m),(1-\rho)A_r(\tfrac12),(1-\rho)A_r(M))$, strictly increasing in the order.
\item In $\mathcal F$, $P(0)-\bar P=-(1-\rho)J_r\,d\,A_r(M)$. The posterior after a neutral order is $\tfrac12$ for every $r$. So the price falls after a neutral order without bad class news if and only if $r\le r_C$.
\item Take the sample of preparation announcements alone. Its mean price change is $\sum_x\Pr(X=x\mid A)\Delta^1_x$, and the mean of the ratios is $\sum_x\Pr(X=x\mid A)\Delta^1_x/P(x)$. Both means are positive in $\mathcal N$, $\mathcal U$ and $\mathcal F$.
\end{enumerate}
The formulas apply to a public event that reveals only preparation.
\end{proposition}

\begin{proof}
(i) After the announcement the dealer knows $x$ and $a\in\{0,1\}$.
Preparation (a=1) gives $\E[V_T]=t_j$ in class $j$.
Since the cost is independent of the class, $\Pr(H\mid x,a=1)=\mu_x$.
So $\E[V_T\mid x,a=1]=t_L+\Delta_T\mu_x=t_0+A_r(\mu_x)$.
Without preparation $\E[V_T\mid x,a=0]=t_0$.
Subtract $P(x)=t_0+e_xA_r(\mu_x)$ (Lemma~\ref{l:prices}).
(ii) Compute $e_x(1-e_x)A-(1-e_x)e_xA=0$.
Divide by the fixed price $P(x)>0$ to get the statement for ratios.
(iii) Immediate from (i), because $A_r(\mu_x)>0$.
The event $A^c$ has probability $1-e_x$ given $X=x$.
(iv) Insert the entries of Table~\ref{tab:tuples}.
(v) By Proposition~\ref{p:spread}(i) and the symmetry of $\Pr(X=\pm1)$,
$\bar P=t_0+\rho A_r(\tfrac12)+(1-\rho)J_rdA_r(M)$.
Here $dA_r(m)+dA_r(M)=2dA_r(\tfrac12)$, since $A_r$ is linear and $m+M=1$.
Subtract this from $P(0)$.
(vi) By (i), $\Delta^1_x\ge0$, with $\Delta^1_x>0$ if and only if $e_x<1$.
The weights $\Pr(X=x\mid A)=\Pr(X=x)e_x/\Ent$ are positive, because $e_x\ge\rho>0$.
In $\mathcal N$, $\mathcal U$ and $\mathcal F$, the orders $x\ne1$ have positive probability and $e_x=\rho<1$.
So both means are positive.
\end{proof}

\paragraph{At B.}
Table~\ref{tab:event} lists the exact changes.
After a buy at $r_1$ the challenger prepares with probability one, so the announcement carries no news.
The no-preparation change after a buy at $r_1$ has probability zero and the table omits it.

\begin{table}[htbp]
\centering\small
\caption{Announcement changes at B: dollars and ratio to $P(x)$.}\label{tab:event}
\begin{tabular}{@{}llcccc@{}}
\toprule
strength & order & $P(x)$ & prepared $\Delta^1$ & ratio & no preparation $\Delta^0$ (ratio)\\
\midrule
$r_1$ & sale & \reg{price_r1_sale_B} & \reg{evA_r1_sale_B} & \reg{evAr_r1_sale_B} & \reg{evN_r1_sale_B} (\reg{evNr_r1_sale_B})\\
$r_1$ & neutral & \reg{price_r1_neutral_B} & \reg{evA_r1_neutral_B} & \reg{evAr_r1_neutral_B} & \reg{evN_r1_neutral_B} (\reg{evNr_r1_neutral_B})\\
$r_1$ & buy & \reg{price_r1_buy_B} & \reg{evA_r1_buy_B} & --- & ---\\
$r_2$ & sale & \reg{price_r2_sale_B} & \reg{evA_r2_sale_B} & \reg{evAr_r2_sale_B} & \reg{evN_r2_sale_B} (\reg{evNr_r2_sale_B})\\
$r_2$ & neutral & \reg{price_r2_neutral_B} & \reg{evA_r2_neutral_B} & \reg{evAr_r2_neutral_B} & \reg{evN_r2_neutral_B} (\reg{evNr_r2_neutral_B})\\
$r_2$ & buy & \reg{price_r2_buy_B} & \reg{evA_r2_buy_B} & \reg{evAr_r2_buy_B} & \reg{evN_r2_buy_B} (\reg{evNr_r2_buy_B})\\
\bottomrule
\end{tabular}
\end{table}

The pre-slot price at $r_1$ is \reg{pbar_r1_B}.
The price after a neutral order is \reg{p_neutral_r1_B}, so the neutral order changes the price by \reg{pdrop_neutral_r1_B}.
At $r_2$ the change is \reg{pdrop_neutral_r2_B}.
Ratios exceed one at $r_2$ because the model sets the standalone value to zero.
A positive standalone value divides the same dollar changes by a larger price.
The mean over both announcements is zero.
The mean over preparation announcements alone is positive.
At $r_1$ it is \reg{evprep_mean_r1_B} in dollars and \reg{evprep_meanr_r1_B} as a return.
At $r_2$ it is \reg{evprep_mean_r2_B} in dollars and \reg{evprep_meanr_r2_B} as a return.

\paragraph{Observable.}
The announcement is the public start of diligence by the further buyer, or a public statement that it will not bid.
The event must reveal only preparation.
The event window starts after the first order.
The mean return given the order is zero over both announcements, with their model probabilities.
A sample of preparation announcements alone has a positive mean.
So a test of zero drift needs the announcements of no preparation.
The model predicts the sign pattern of the surprises across orders and across the ceiling.

\section{The coexistence branch}\label{sec:coex}

\begin{proposition}[Coexistence; analytical, values computer-assisted]\label{p:coex}
Let $r_F<r<r_N$ and $r\le r_C(M)$, so that $\mathcal N$ and $\mathcal U$ are both equilibrium outcomes (Proposition~\ref{p:out}(c)).
\begin{enumerate}[label=(\roman*)]
\item \emph{Order flow.} $\mathcal N$: $\E[X]=0$ and $v=2b$. $\mathcal U$: $\E[X]=\lambda/2>0$ and $v=2b+\lambda/2$.
\item \emph{Prices.} In $\mathcal N$, $P(-1)=P(0)=P(1)$. In $\mathcal U$, $P(0)<P(-1)<P(1)$, with
\[
P(-1)-P(0)=\rho\Delta_T\,\frac{\lambda}{4(1-v_{\mathcal U})},\qquad v_{\mathcal U}=2b+\tfrac\lambda2 .
\]
\item \emph{Entry.} $\Pr(A\mid X=1)=1$ in $\mathcal U$. All other entry rates equal $\rho$ in both outcomes. $\Ent=\rho$ in $\mathcal N$ and $\Ent=\rho+(1-\rho)d$ in $\mathcal U$.
\item \emph{Association.} In $\mathcal N$ the buy impact $P(1)-P(0)$ is zero and the entry after a buy is $\rho$, so the entry increment $\Pr(A\mid X=1)-\rho$ is zero. In $\mathcal U$ the buy impact is positive and the entry increment is $1-\rho$. Hence, across the two outcomes at one public strength, the buy impact is positive if and only if the entry increment is positive.
\item \emph{Classification.} Under the challenger rule (an indifferent challenger prepares), $\mathcal N$ and $\mathcal U$ are the only equilibrium outcomes. The investor rule does not matter, because no investor is indifferent.
\end{enumerate}
\end{proposition}

\begin{proof}
(i) is Proposition~\ref{p:flow}.
(ii) Table~\ref{tab:tuples} gives $\mu_{-1}=\tfrac12$ and $\mu_0^{\mathcal U}=(1-2b-\lambda)/(2-4b-\lambda)$.
Then $\tfrac12-\mu_0^{\mathcal U}=\lambda/[2(2-4b-\lambda)]$.
Since $4b+\lambda=2v_{\mathcal U}$, this equals $\lambda/[4(1-v_{\mathcal U})]>0$.
Both orders have entry $\rho$, so $P(-1)-P(0)=\rho\Delta_T(\tfrac12-\mu_0^{\mathcal U})$.
The price $P(1)$ exceeds both because $e_1=1$ and $A_r(M)>A_r(\tfrac12)$.
(iii) Table~\ref{tab:tuples} with $J_r=1$.
(iv) In $\mathcal U$, $P(1)-P(0)>0$ and entry after a buy is one, so the entry increment is $1-\rho>0$.
In $\mathcal N$, the buy impact is zero and the entry after a buy is $\rho$.
So the entry increment above $\rho$ is zero.

(v) Take any equilibrium.
A high-class sale loses, $U_H(-1)=-e_{-1}\Delta_T(1-\mu_{-1})-k<0$, and a low-class buy loses, $U_L(1)=-e_1\Delta_T\mu_1-k<0$.
So $\sigma_H(-1)=0$ and $\sigma_L(1)=0$.
Let $t=\lambda\sigma_L(-1)$ and $s=\lambda\sigma_H(1)$.
Then $\mu_{-1}=b/(2b+t)\le\tfrac12$, so $B_r(\mu_{-1})\le B_r(\tfrac12)<\cH$ and the entry after a sale is $\rho$.
The low-class sale payoff is $\rho\Delta_T\mu_{-1}-k\le\rho\Delta_T/2-k<0$, because $r<r_N$.
So $t=0$.
Then $\mu_0=(1-2b-s)/(2-4b-s)\le\tfrac12$, so $B_r(\mu_0)<\cH$ and the neutral order does not recruit.
The buy posterior is $\mu_1=(b+s)/(2b+s)\in[\tfrac12,M]$, and the challenger rule gives $e_1\in\{\rho,1\}$.
If $e_1=\rho$, then $U_H(1)=\rho\Delta_T(1-\mu_1)-k\le\rho\Delta_T/2-k<0$.
So $s=0$, and this is $\mathcal N$.
If $e_1=1$, then $U_H(1)=\Delta_T(1-\mu_1)-k\ge\Delta_T m-k>0$, because $\mu_1\le M$ and $r>r_F$.
So $s=\lambda$, and this is $\mathcal U$.
\end{proof}

\begin{remark}[Partial buys when an indifferent challenger may mix]\label{r:mix}
The classification (v) depends on the challenger rule.
Let $r_F<r<r_N$ and $r<r_C(M)$.
Write $\tau=\tau(r)$, so that $B_r(\tau)=\cH$ and $\tau>\tfrac12$.
Let the challenger with cost $\cH$ prepare with probability $e_1\in[\rho,1]$ at the posterior $\tau$.
Set
\[
s_\tau=b\,\frac{2\tau-1}{1-\tau},\qquad t=0,\qquad \sigma_H(1)=\frac{s_\tau}{\lambda},\qquad e_1=\frac{k}{\Delta_T(1-\tau)} .
\]
Then $\mu_1=(b+s_\tau)/(2b+s_\tau)=\tau$, so the challenger is indifferent.
The probability $\sigma_H(1)$ lies in $(0,1)$, because $\tfrac12<\tau<M$.
The high class earns $e_1\Delta_T(1-\tau)-k=0$ from a buy, so it is indifferent between a buy and no trade.
The probability $e_1$ lies in $(\rho,1)$: $\Delta_T(1-\tau)>\Delta_T m>k$ gives $e_1<1$, and $\rho\Delta_T(1-\tau)<\rho\Delta_T/2<k$ gives $e_1>\rho$.
The low class earns $\rho\Delta_T/2-k<0$ from a sale, so it strictly abstains.
All non-buy entry rates equal $\rho$, and the price reveals the order.
The tuple is an equilibrium outcome under the declared investor rule and a challenger rule that allows a mixture at indifference.
It adds a third outcome to Proposition~\ref{p:coex}(v).
At $r=1.60$ the high class buys with probability \reg{mix_sigma_r160_B}, and the entry after a buy is \reg{mix_e1_r160_B}.
The mean order is \reg{mix_mean_order_r160_B}, the entry is \reg{mix_entry_r160_B}, and the low-class sale payoff is \reg{mix_lowsale_r160_B}.
\end{remark}

\paragraph{At B.}
At $1.55$, $1.60$, $1.65$ the buy-only branch exists with margins \reg{margin_obs_rF_B} above $r_F=\reg{rF_check_B}$ and \reg{margin_obs_rN_B} below $r_N=\reg{rN_check_B}$.
The spread $P(1)-P(-1)$ is \reg{ba_U_r155_B}, \reg{ba_U_r160_B}, \reg{ba_U_r165_B}.
At $1.60$ a sale executes \reg{U_gap_r160_B} above the no-trade price.
The posterior after a neutral order is \reg{U_mu0_B}.
The mean order is \reg{imb_buy_B}.
Entry is \reg{U_entry_B}, the same as at $r_1$.
The ceiling lies \reg{margin_obs_rN_rC_B} above $r_N$, so $\mathcal U$ recruits on the whole interval.

\paragraph{Observable.}
The branch separates by the sign of the mean order, by the sign of $P(-1)-P(0)$, and by the entry rate after a buy.
Under the challenger rule, a cross-section of processes with strength in $(r_F,r_N)$ shows two clusters: zero price impact with entry $\rho$, and positive price impact with entry after a buy equal to one.
If an indifferent challenger may mix, a third cluster appears (Remark~\ref{r:mix}).
It has a mean order between $0$ and $\lambda/2$ and an entry after a buy between $\rho$ and $1$.

\section{Comparison with the earlier study}\label{sec:old}

The earlier study (memo N1) derived four extensions at $\beta=1/3$ and $\lambda=6/17$ (calibration A).
Section~10 of the run log recomputes 21 of its exact values in rational arithmetic.
All match.
Table~\ref{tab:old} states what survives at B.
One structural statement of the earlier study fails.
It says that each spread component rises strictly with the strength while $J$ stays fixed.
On the regime $J=0$ the entry component is identically zero (Proposition~\ref{p:spread}(v)).
Every number that depends on $(\lambda,\beta)$ changes.

\begin{table}[htbp]
\centering\small
\caption{Claims of the earlier study at B.}\label{tab:old}
\begin{tabular}{@{}p{3.8cm}p{3.2cm}p{3.3cm}p{3.3cm}@{}}
\toprule
Claim & Earlier value & Value at B & Verdict\\
\midrule
Spread $=$ class $+$ entry component & $S(3)=25/32$ & \reg{ba_r1_B} & identity survives; value changes\\
Entry share of the spread at $r_1$ & $90.4\%$ & \reg{ba_share_r1_B} & value changes\\
Spread falls at $r_C$ & $r_C=3.458247$ & $r_C=\reg{rC_check_B}$; fall \reg{ba_jump_rC_B} & survives\\
Each spread component rises strictly with $r$ while $J$ is fixed & stated & class component rises; entry component is $0$ for $J=0$ & fails on $J=0$; corrected\\
Informed share on $\mathcal F$ & $9/20$ & \reg{theta_B} & value changes\\
Informed share and volume constant from $r_1$ to $r_2$ & yes & yes & survives\\
Identification from sell impact and volume & recovers $(9/20,6/17,1/3)$ & recovers $(\vartheta,\lambda,\beta)$ & survives\\
Counts do not identify the share & pairs $(6/17,1/3)$, $(1/3,23/68)$ & pair $u'=8/5$ and the range $[\reg{theta_min_counts_B},\reg{theta_max_full_eq_r1_B})$ & survives; the pair with $\beta'=23/68$ lies outside the model ($\beta'>1/3$)\\
Buy-only branch: volume, share, mean order & $31/51$, $9/31$, $3/17$ & \reg{vol_buy_B}, \reg{theta_buy_B}, \reg{imb_buy_B} & values change\\
Winner gap $\vartheta(1-\ell/r)$ & $3/10$, $13/40$ & \reg{wgap_r1_B}, \reg{wgap_r2_B} & survives; values change\\
Entry and winner covariances & $5/17$, $13/408$ & \reg{cov_entry_r1_B}, \reg{cov_win_r2_B} & survive\\
Neutral order lowers the price & $-113/408$ & \reg{pdrop_neutral_r1_B} & survives\\
Announcement changes & $(77/160,19/32,0)$ at $r_1$ & Table~\ref{tab:event} & survive; values change\\
Full orders unique on $[3,18/5]$ & yes & unique on $(r_S,h)$, $r_S=\reg{rS_p5_B}$ (rounded) & extended\\
Informed share $=2M-1$ & ``appropriate restrictions'' & holds on $\mathcal F$ at $\pi=\tfrac12$ & stated exactly\\
\bottomrule
\end{tabular}
\end{table}

The informed share of nonzero orders on $\mathcal F$ is $\lambda/(\lambda+2b)=u/(2+u)$ at every prior.
The posterior gap $M(\pi)-m(\pi)$ equals this number only at $\pi=\tfrac12$.
The earlier study noted the restriction.
Here it is Proposition~\ref{p:flow}.

\section{Margins, tie rules, and the parameter set}\label{sec:margins}

\subsection{Margins at B}

Every hypothesis is a strict inequality.
Table~\ref{tab:margins} lists the margins.
The smallest margin is the return to the entry floor at $r_2$, which is \reg{margin_obs_ceiling_B} in posterior.
The primitive margins are positive: $\lambda=\reg{lambda_p5_B}$, $1-\lambda$, $\beta=\reg{beta_p5_B}$, $\rho=1/4$, $1-\rho=3/4$ and $k=1/50$.
At the three benchmark strengths, the smallest $r-\ell$ is $1/5$ and the smallest $h-r$ is $32/5$.

\begin{table}[htbp]
\centering\small
\caption{Hypotheses of the package and their margins at B.}\label{tab:margins}
\begin{tabular}{@{}p{6.4cm}p{3.2cm}p{3.5cm}@{}}
\toprule
Hypothesis & Margin & Used by\\
\midrule
$\beta<1/3$ & $1/3-\beta=\reg{margin_obs_beta_B}$ & Lemma~\ref{l:bounds}\\
$\Delta_T(r_0)<k$ & $k-\Delta_T(r_0)=\reg{margin_obs_weak_r0_B}$ & Prop.~\ref{p:out}(a)\\
$\rho\Delta_T(r_1)m>k$ & \reg{margin_obs_full_orders_r1_B} & Prop.~\ref{p:out}(b)\\
$\rho\Delta_T(r_2)m>k$ & \reg{margin_obs_full_orders_r2_B} & Prop.~\ref{p:out}(b)\\
$M>\tau(r_1)$ (buy recruits at $r_1$) & \reg{margin_obs_buy_recruit_B} (posterior) & $J_{r_1}=1$\\
$M<\tau(r_2)$ (no recruit at $r_2$) & \reg{margin_obs_ceiling_B} (posterior) & $J_{r_2}=0$\\
$M>5/8$ (the ceiling exists) & \reg{margin_obs_post_over_58_B} (posterior) & Lemma~\ref{l:auction}(d)\\
$r_S<r_C$ (the ceiling lies in the $\mathcal F$ interval) & $r_C-r_S=\reg{margin_obs_rS_rC_B}$ & Prop.~\ref{p:spread}(v), (vi)\\
$r_N<r_C$ (recruitment on the coexistence interval) & $r_C-r_N=\reg{margin_obs_rN_rC_B}$ & Prop.~\ref{p:coex}\\
$r_S<h$ (the $\mathcal F$ interval is not empty) & $h-r_S=\reg{margin_obs_rS_h_B}$ & Prop.~\ref{p:out}(b)\\
$r_C<r_2$ & $r_2-r_C=\reg{margin_obs_r2_above_rC_B}$ & Prop.~\ref{p:spread}(vi)\\
$B_{r_1}(\tfrac12)<\cH$ & $\cH-B=\reg{margin_obs_neutral_r1_B}$ & neutral order\\
$u>82/59$ (lower band) & \reg{margin_obs_u_lo_B} & band\\
$2686/1559-u>0$ (upper band) & \reg{margin_obs_u_hi_B} & band\\
$\rho\Delta_T(r_1)(1-\tau(r_2))>k$ (incentives on the band) & \reg{margin_obs_band_incentive_B} & Prop.~\ref{p:ident}(c), (d)\\
$\Delta_T(1.55)m>k$ (buy-only) & \reg{margin_obs_buy_r155_B} & Prop.~\ref{p:out}(c)\\
$B_{1.65}(M)>\cH$ (buy recruits) & \reg{margin_obs_buy_recruit_r165_B} & Prop.~\ref{p:out}(c)\\
$k-\rho\Delta_T(1.65)/2>0$ (low class waits) & \reg{margin_obs_lowsell_r165_B} & Prop.~\ref{p:out}(c)\\
$1.55>r_F$ & \reg{margin_obs_rF_B} & Prop.~\ref{p:coex}\\
$r_N>1.65$ & \reg{margin_obs_rN_B} & Prop.~\ref{p:coex}\\
\bottomrule
\end{tabular}
\end{table}

\subsection{Tie rules}

There are two tie rules.
The \emph{investor rule} lets an indifferent investor choose any mixture.
The \emph{challenger rule} makes an indifferent challenger prepare.
Table~\ref{tab:ties} shows which results depend on them.

\begin{table}[htbp]
\centering\small
\caption{Dependence on the tie rules.}\label{tab:ties}
\begin{tabular}{@{}p{5.2cm}p{3.0cm}p{4.9cm}@{}}
\toprule
Result & Investor rule & Challenger rule\\
\midrule
Props.~\ref{p:flow}, \ref{p:ident}, \ref{p:win}, \ref{p:event} at $r_1$, $r_2$ & none (strict margins) & none ($r_1,r_2\ne r_C$)\\
Prop.~\ref{p:spread} and Prop.~\ref{p:entry} at $r\ne r_C$ & none & none\\
Spread, buy impact, $\bar P$ and entry exactly at $r_C$ & none & prepare: $J=1$. With a mixture, each lies between its two one-sided limits.\\
Prop.~\ref{p:coex}(i)--(iv) at $r_F<r<r_N$: existence and formulas of $\mathcal N$ and $\mathcal U$ & none (strict) & none\\
Prop.~\ref{p:coex}(v): $\mathcal N$ and $\mathcal U$ are the only outcomes & none (no investor is indifferent) & needs the rule \emph{prepare}. A mixture at indifference adds the partial-buy outcome (Remark~\ref{r:mix}).\\
Boundaries $r_F$, $r_N$, $r_S$ & the outcome at the boundary needs the rule & none\\
\bottomrule
\end{tabular}
\end{table}

\subsection{The parameter set}

The hypotheses define the set $\Theta$ of Table~\ref{tab:theta} in the two free slot parameters $(\lambda,\beta)$, at the prior $1/2$.
Each restriction is strict, so $\Theta$ is open in the plane.
The set contains B.

\begin{table}[H]
\centering\small
\caption{The open parameter set $\Theta$.}\label{tab:theta}
\begin{tabular}{@{}p{4.0cm}p{8.9cm}@{}}
\toprule
Free parameter & Strict restriction\\
\midrule
$\beta$ & $0<\beta<1/3$\\
$\lambda$ & $\dfrac{(82/59)\beta}{1+(82/59)\beta}<\lambda<\dfrac{(2686/1559)\beta}{1+(2686/1559)\beta}$\\
Equivalent posterior & $141/200<M<4245/5804$\\
Strength in $\mathcal F$ & $r_S(M)<r<h$\\
Strength in the coexistence interval & $r_F(M)<r<r_N$\\
\bottomrule
\end{tabular}
\end{table}

The band for $\lambda$ is the band $82/59<u<2686/1559$ for $u$, because $\lambda=u\beta/(1+u\beta)$ increases in $u$.
It is the band $\tau(r_1)<M<\tau(r_2)$ with $\tau(r_1)=141/200$ and $\tau(r_2)=4245/5804$.
On $\Theta$ the following statements hold.
\begin{enumerate}[label=(\alph*)]
\item $M>141/200>5/8$, so the ceiling $r_C(M)$ exists.
\item $r_C(M)$ increases in $M$ and $B_{r_1}(M)>\cH>B_{r_2}(M)$, so $3<r_C(M)<3.6$.
\item $\rho\Delta_T(r_1)m>\rho\Delta_T(r_1)(1-\tau(r_2))>k$, so $r_S(M)<3$.
\item $r_F(M)<r_N<r_S(M)$, because $m>1/8$ gives the first inequality and $m<1/2$ gives the second.
\item $r_N<3<r_C(M)$, so the coexistence interval lies below the ceiling.
\end{enumerate}
The uniqueness margin $\rho\Delta_T(r)m-k$ tends to zero as $r\downarrow r_S$.
So no uniform positive margin holds on the whole interval $(r_S,h)$.

The cases $\beta=1/3$, $r=r_C$ and $r=r_N$ are boundary cases and lie outside $\Theta$.
At $\beta=1/3$ the bound of Lemma~\ref{l:bounds} holds with equality at the neutral order, and every proof uses only $\beta\le1/3$.
The results therefore hold at $\beta=1/3$.
At $r=r_C$ the challenger rule decides the entry after a buy.
At $r=r_N$ and at $r=r_S$ the low class is indifferent, and the investor rule decides the outcome.

\subsection{Reproduction}

Run from the repository root:
\begin{verbatim}
.venv/bin/python audit/slot_manuscript_2026_10_23/scripts/P5_observables/p5_observables.py \
  > audit/slot_manuscript_2026_10_23/scripts/P5_observables/run_log.txt
\end{verbatim}
The run takes less than one second.
It writes the registry rows, the file \texttt{p5\_numbers.tex} that this document reads, and \texttt{curves\_full\_branch.csv} (figure data on $2.2\le r\le6$).

\end{document}
